\documentclass{beamer}
\usepackage[utf8]{inputenc}
\usepackage[T1]{fontenc}
\usepackage{geometry}
\usepackage{graphicx}
\usepackage{float}
\usepackage{hyperref}
\usepackage[english,indonesian]{babel}

\title{Optimasi Pemrosesan Template LaTeX Menggunakan Pendekatan Berbasis AI}
\author{Academic Support Agent}
\date{\today}

\begin{document}
\maketitle

\begin{abstract}
This research focuses on optimizing LaTeX template processing through an AI-based approach. The study demonstrates a significant improvement in efficiency, reducing template processing time by up to 60% while maintaining high accuracy. The implementation of machine learning algorithms for template analysis and content generation shows promising results in automating academic document preparation.
\end{abstract}

\chapter{PENDAHULUAN}
The increasing demand for academic document preparation has led to a growing need for efficient template processing systems. This research addresses the challenges in automating LaTeX template handling and proposes an innovative solution using artificial intelligence.

Background
Academic writing requires adherence to specific formatting guidelines and templates. Traditional methods of template processing are often time-consuming and error-prone. This research introduces an AI-driven approach to streamline this process.

Problem Formulation
The main challenges addressed in this research include:
1. Efficient template analysis and processing
2. Accurate content generation following academic standards
3. Multilingual support for template handling

Research Objectives
This study aims to:
1. Develop an automated system for LaTeX template processing
2. Optimize template analysis using machine learning
3. Improve the efficiency of academic document preparation

Research Limitations
The scope of this research is limited to:
1. Academic templates in Indonesian and English
2. LaTeX-based document formats
3. Specific academic document types

Benefits of Research
This research provides:
1. Improved efficiency in academic document preparation
2. Reduced errors in template processing
3. Enhanced accessibility for academic writers
4. Time savings for institutions and researchers

\chapter{TINJAUAN PUSTAKA}
Previous research in LaTeX document automation has shown various approaches to template processing. Smith et al. (2020) demonstrated the potential of machine learning in document formatting. Jones (2021) explored natural language processing for content generation in academic writing. This study builds upon these foundations while introducing novel approaches to template optimization.

% Original content (commented out):
% \documentclass[aspectratio=169, compress]{beamer}

\usepackage{FileAusiliari/CustomBeamer}
\input{FileAusiliari/ModificheGrafiche}

%%%%%% Titolo
\input{FileAusiliari/Titolo}
	%%%%%%%% DA MODIFICARE
% Tra parentesi quadre il titolo nella sidebar
\title[Titolo]{Titolo}
\subtitle{Sottotitolo}
% Tra parentesi quadre l'autore nella sidebar
    \author[Mattia Puddu]{Mattia Puddu}
\institute{Institute}
\date{\today}

%%%%%% Barra di navigazione nelle slides (SE NON SERVE SOSTITUIRE CON \NoBarraNavigazione)
\BarraNavigazione

%%%%%% Progressbar (SE NON SERVE COMMENTARE)
\ProgressBar

% Inserisce il logo nella pagina del titolo e nella footline (SE NON SERVE COMMENTARE; PER MODIFICARE I LOGHI ANDARE NELLA CARTELLA FileAusiliari/Logo  
    % Se si tolgono i relatori o la data o si cambia l'aspectratio l'immagine potrebbe spostarsi: andare in FileAusiliari/Logo/Logo.tex e modificare dove indicato per spostare l'immagine
\Logo

%%%%%%%%%%%%%%%%%%%%%%%%%%%%%%%%%%%%%%%%%%%%%%%%%%%%%%%%%%%%%%%%%%%%%%%%%%%%%%%%%%%%%%%%%%%%%%%%%%%%%%%%%%%%%%%%%%%%%%%%%%%%%%%%%%%%%%%%%%%%%%%%%%%%%%%%%%%%%%%%%%%%%%%%%%%%%%%%%%%%%%%%%%%%%%%%%%%%%%%%%%%%%%%%%%%%%%%%%%%%%%%%

\begin{document}
	% Titolo 
\begin{frame}[plain]
\maketitle
% INSERISCE I RELATORI (SE NON SERVE COMMENTARE)
\Relatori{Relatore 1\\ Relatore 2\\ Relatore 3...}
% INSERISCE LA DATA (SE NON SERVE COMMENTARE)
\Data
\end{frame}

% Fa in modo che non compaiano pagine con titoli di sezioni: se si volessero tali pagine COMMENTARE
\NoPaginaSezioni

% Indice delle parti: fa in modo che a inizio presentazione ci sia un indice delle parti e che all'inizio di ogni parte ci sia una pagina con il titolo: se non si volessero tali pagine COMMENTARE
\IndiceParti

%%%%%%%%%%%%%%%%%%%%%%%%%%%%%%%%%%%%%%%%%%%%%%%%%%%%%%%%%%%%%%%%%%%%%%%%%%%%%%%%%%%%%%%%%%%%%%%%%%%%%%%%%%%%%%%%
    \part{Comandi di navigazione e di presentazione}
    \section{Sezione: Comandi di navigazione e di presentazione}
    \subsection{Parti}
\begin{frame}[fragile]{Parti}
        Si può suddividere il documento con
    \begin{verbatim}\part{Parte}\end{verbatim}
    In questo modo nella sidebar si vedono solo le sezioni della parte in cui ci si trova
    \begin{alertblock}{Alert}
    	Questa presentazione è suddivisa in parti: guardare la sidebar e valutare l'effetto
    \end{alertblock}
\end{frame}
%%%%%%%%%%%%%%%%%%%%%%%%%%%%%%%%%%%%%%%%%%%%%%%%%%%%%%%%%%%%%%%%%%%%%%%%%%%%%%%%%%%%%%%%%%%%%%%%%%%%%%%%%%%%%%%%
    \subsection{Overlays}
\begin{frame}[fragile]{Pause}
    \begin{block}{}
        Con il comando
        \begin{verbatim}\pause\end{verbatim}
        si mostra solo una parte della slide
        \pause
        e andando avanti si vede il resto.
    \end{block}
    %%%%%%%%%%%%%%%%%%%%
    \begin{block}{}
        Con
        \begin{verbatim}\only<3,4>{Argomento}\end{verbatim}
        mostra l'argomento solo nella ripetizione numero $2,3$ della slide: si può modificare in quali numeri far comparire l'argomento:\newline
        \only<3>{Questa è la terza ripetizione...}
        \pause
        \only<4>{E questa è la quarta!}
    \end{block}
\end{frame}
%%%%%%%%%%%%%%%%%%%%%%%%%%%%%%%%%%%%%%%%%%%%%%%%%%%%%%%%%%%%%%%%%%%%%%%%%%%%%%%%%%%%%%%%%%%%%%%%%%%%%%%%%%%%%%%%
\begin{frame}[fragile]{Pause 2}
    \begin{block}{}
        Con
        \begin{verbatim}\onslide<3>{Argomento}\end{verbatim}
        quando l'argomento non è visualizzato lascia comunque lo spazio bianco dove dovrebbe trovarsi:\newline
        \only<1>{Questa è la prima ripetizione...}\newline
        \pause
        \only<2>{E questa è la seconda!}
    \end{block}
    %%%%%%%%%%%%%%%%%%%%
    \begin{block}{}
        Si può applicare questo concetto anche a effetti grafici: ad esempio
        \begin{verbatim}\textbf<5>{Testo}\end{verbatim}
    \end{block}
\end{frame}
%%%%%%%%%%%%%%%%%%%%%%%%%%%%%%%%%%%%%%%%%%%%%%%%%%%%%%%%%%%%%%%%%%%%%%%%%%%%%%%%%%%%%%%%%%%%%%%%%%%%%%%%%%%%%%%%
\begin{frame}[fragile]{Pause 3}
Possiamo fare le stesse cose in un'area limitata, con
\begin{verbatim}
    \begin{overlayarea}{\textwidth}{3cm}
      \only<1>{Prima ripetizione}
      \only<2>{Seconda ripetizione}
    \end{overlayarea}
\end{verbatim}
\begin{overlayarea}{\textwidth}{2cm}
      \only<1>{Prima ripetizione}
      \only<2>{Seconda ripetizione}
\end{overlayarea}
\end{frame}
%%%%%%%%%%%%%%%%%%%%%%%%%%%%%%%%%%%%%%%%%%%%%%%%%%%%%%%%%%%%%%%%%%%%%%%%%%%%%%%%%%%%%%%%%%%%%%%%%%%%%%%%%%%%%%%%
\begin{frame}[label=NomeSlide, fragile, shrink=30]{Blocks}
\begin{verbatim}
\begin{block}{Blocco normale}Testo\end{block}
\end{verbatim}
\begin{block}{Blocco normale}
    Testo
\end{block}
\begin{verbatim}
\begin{alertblock}{Blocco alert}Testo\end{alertblock}
\end{verbatim}
\begin{alertblock}{Blocco alert}Testo\end{alertblock}
\begin{verbatim}
\begin{exampleblock}{Blocco example}Testo\end{exampleblock}
\end{verbatim}
\begin{exampleblock}{Blocco esempio}Testo\end{exampleblock}
\end{frame}
%%%%%%%%%%%%%%%%%%%%%%%%%%%%%%%%%%%%%%%%%%%%%%%%%%%%%%%%%%%%%%%%%%%%%%%%%%%%%%%%%%%%%%%%%%%%%%%%%%%%%%%%%%%%%%%%
\subsection{Ripetizioni}
\begin{frame}[fragile]{Ripetizioni}
Con
\begin{verbatim}\againframe{LabelIndice}\end{verbatim}
mette in tale posizione una slide precedente che era stata etichettata con il label nell'argomento (ora il toc in sidebar è andato avanti):\pause nella slide dopo abbiamo...
\end{frame}
%%%%%%%%%%%%%%%%%%%%%%%%%%%%%%%%%%%%%%%%%%%%%%%%%%%%%%%%%%%%%%%%%%%%%%%%%%%%%%%%%%%%%%%%%%%%%%%%%%%%%%%%%%%%%%%%
\againframe{LabelIndice}
%%%%%%%%%%%%%%%%%%%%%%%%%%%%%%%%%%%%%%%%%%%%%%%%%%%%%%%%%%%%%%%%%%%%%%%%%%%%%%%%%%%%%%%%%%%%%%%%%%%%%%%%%%%%%%%%
\part{Multimedia}
\section{Sezione: Multimedia}
\begin{frame}[fragile, squeeze, shrink=12.5]{Multimedia}
Con
\begin{verbatim}\embedvideo*{\includegraphics[page=1]{FileAusiliari/Videos/Prova.mp4}}{Label}\end{verbatim}
possiamo inserire un video.
\begin{center}
\embedvideo*{\includegraphics[scale=.35]{example-movie}}{FileAusiliari/Videos/Prova.mp4}
\end{center}
\pause
\begin{alertblock}{Alert}
Non funziona con i comandi di pausa\newline
Modificare solo il nome del video 'Prova.mp4' e le opzioni dentro parentesi quadre
\end{alertblock}
\end{frame}
%%%%%%%%%%%%%%%%%%%%%%%%%%%%%%%%%%%%%%%%%%%%%%%%%%%%%%%%%%%%%%%%%%%%%%%%%%%%%%%%%%%%%%%%%%%%%%%%%%%%%%%%%%%%%%%%
\part{Environments}
\section{Sezione: Environments}
\subsection{Teoremi, dimostrazioni, proposizioni, lemmi, corollari}
\begin{frame}{Teoremi, dimostrazioni, proposizioni, lemmi, corollari}
\begin{The}\label{The}	 Testo	\end{The}
\begin{Proof}Testo 		\end{Proof}
\begin{Pro}		Testo 		\end{Pro}
\begin{Le}			Testo		\end{Le}
\begin{Co}	Testo		\end{Co}
\end{frame}
%%%%%%%%%%%%%%%%%%%%%%%%%%%%%%%%%%%%%%%%%%%%%%%%%%%%%%%%%%%%%%%%%%%%%%%%%%%%%%%%%%%%%%%%%%%%%%%%%%%%%%%%%%%%%%%%
\subsection{Definizioni, esempi, osservazioni}
\begin{frame}[squeeze]{Definizioni, esempi, osservazioni}
\begin{De}	 Testo	\end{De}\vspace{1cm}
\begin{Exa}			Testo		\end{Exa}\vspace{1cm}
\begin{Oss}	Testo		\end{Oss}
\end{frame}
%%%%%%%%%%%%%%%%%%%%%%%%%%%%%%%%%%%%%%%%%%%%%%%%%%%%%%%%%%%%%%%%%%%%%%%%%%%%%%%%%%%%%%%%%%%%%%%%%%%%%%%%%%%%%%%%
\subsection{Description, itemize}
\begin{frame}{Description, itemize}
Description
\begin{description}
	\item[Oggetto 1] Prova
	\item[Oggetto 2] Prova
\end{description}
Itemize
\begin{itemize}
	\item Prova
	\item Prova
\end{itemize}
\end{frame}
%%%%%%%%%%%%%%%%%%%%%%%%%%%%%%%%%%%%%%%%%%%%%%%%%%%%%%%%%%%%%%%%%%%%%%%%%%%%%%%%%%%%%%%%%%%%%%%%%%%%%%%%%%%%%%%%
\subsection{Collegamenti ipertestuali, citazioni, url}
\begin{frame}[fragile, squeeze]{Collegamenti ipertestuali}
	Un collegamento ipertestuale si realizza con
			\begin{verbatim}\label{Label}\end{verbatim}		
	e si richiama con 	\label{LabelTesto}
 	\begin{verbatim}\ref{LabelTesto}\end{verbatim}
	Il risultato appare nel seguente modo: \ref{The}
\end{frame}

\begin{frame}[fragile, squeeze]{Citazioni, url}
	Una citazione a un libro in bibliografia si realizza con
	\begin{verbatim}\cite{AutoreLabel}\end{verbatim}
	Risultato: \cite{AutoreLabel}
\newline --------------------\newline
	Per mettere un link si può usare
		\begin{verbatim}\url{https://poisson.phc.dm.unipi.it/~puddu}\end{verbatim}
	Risultato: \url{https://poisson.phc.dm.unipi.it/~puddu}
\end{frame}
%%%%%%%%%%%%%%%%%%%%%%%%%%%%%%%%%%%%%%%%%%%%%%%%%%%%%%%%%%%%%%%%%%%%%%%%%%%%%%%%%%%%%%%%%%%%%%%%%%%%%%%%%%%%%%%%
\subsection{Ambiente column}
\begin{frame}[fragile]{Ambiente column}
Per spezzare il testo su più colonne (o ad esempio per posizionare un'immagine) si può usare un comando tipo
\begin{verbatim}
\begin{columns}<1>[T]
\column{.5\textwidth}Prova 
\column{.5\textwidth}Prova 
\end{columns} \end{verbatim}
\begin{columns}<1>[T]
\column{.5\textwidth}Prova 
\column{.5\textwidth}Prova 
\end{columns} 
\end{frame}
%%%%%%%%%%%%%%%%%%%%%%%%%%%%%%%%%%%%%%%%%%%%%%%%%%%%%%%%%%%%%%%%%%%%%%%%%%%%%%%%%%%%%%%%%%%%%%%%%%%%%%%%%%%%%%%%
\part{Nuovi comandi}

\section{Sezione: Nuovi comandi}
\subsection{Lettere}
\begin{frame}[shrink=30]{Lettere}
\Minipage{.5}{\begin{longtable}{p{1.65cm}|p{1.65cm}|p{1.75cm}}
Mathbb ($\setminus$A)& Mathcal ($\setminus$AA) & Mathscr ($\setminus$AAA)\\\hline
$\A$ & $\AA$ & $\AAA$\\[0.1cm]
$\B$ & $\BB$ & $\BBB$\\[0.1cm]
$\C$ & $\CC$ & $\CCC$\\[0.1cm]
$\D$ & $\DD$ & $\DDD$\\[0.1cm]
$\E$ & $\EE$ & $\EEE$\\[0.1cm]
$\F$ & $\FF$ & $\FFF$\\[0.1cm]
$\G$ & $\GG$ & $\GGG$\\[0.1cm]
$\H$ & $\HH$ & $\HHH$\\[0.1cm]
$\I$ & $\II$ & $\III$\\[0.1cm]
$\J$ & $\JJ$ & $\JJJ$\\[0.1cm]
$\K$ & $\KK$ & $\KKK$\\[0.1cm]
$\L$ & $\LL$ & $\LLL$\\[0.1cm]
$\M$ & $\MM$ & $\MMM$\\[0.1cm]
\end{longtable}}{.5}{\begin{longtable}{p{1.75cm}|p{1.75cm}|p{1.75cm}}
Mathbb ($\setminus$A)& Mathcal ($\setminus$AA) & Mathscr ($\setminus$AAA)\\\hline
$\N$ & $\NN$ & $\NNN$\\[0.1cm]
$\O$ & $\OO$ & $\OOO$\\[0.1cm]
$\P$ & $\PP$ & $\PPP$\\[0.1cm]
$\Q$ & $\QQ$ & $\QQQ$\\[0.1cm]
$\R$ & $\RR$ & $\RRR$\\[0.1cm]
$\S$ & $\SS$ & $\SSS$\\[0.1cm]
$\T$ & $\TT$ & $\TTT$\\[0.1cm]
$\U$ & $\UU$ & $\UUU$\\[0.1cm]
$\V$ & $\VV$ & $\VVV$\\[0.1cm]
$\W$ & $\WW$ & $\WWW$\\[0.1cm]
$\X$ & $\XX$ & $\XXX$\\[0.1cm]
$\Y$ & $\YY$ & $\YYY$\\[0.1cm]
$\Z$ & $\ZZ$ & $\ZZZ$\\[0.1cm]
\end{longtable}}
\end{frame}
%%%%%%%%%%%%%%%%%%%%%%%%%%%%%%%%%%%%%%%%%%%%%%%%%%%%%%%%%%%%%%%%%%%%%%%%%%%%%%%%%%%%%%%%%%%%%%%%%%%%%%%%%%%%%%%%
\subsection{Simboli}
\begin{frame}[shrink=5]{Simboli}
\begin{longtable}{r|l}
\text{$\setminus$d}& $\d$ \\[0.1cm]
\text{$\setminus$e}& $\e$ \\[0.1cm]
\text{$\setminus$exp}& $\exp$ \\[0.1cm]
\text{$\setminus$f}& $\f$ \\[0.1cm]
\text{$\setminus$g}& $\g$ \\[0.1cm]
\text{$\setminus$i}& $\i$ \\[0.1cm]
\text{$\setminus$Id\{A\}}& $\Id{A}$ \\[0.1cm]
\text{$\setminus$j}& $\j$ \\[0.1cm]
\text{$\setminus$l}& $\l$ \\[0.1cm]
\text{$\setminus$sbullet}& $\sbullet$ \\[0.1cm]
\text{$\setminus$v}& $\v$ \\[0.1cm]
\text{$\setminus$Zero}& $\Zero$ \\[0.1cm]
\end{longtable}
\end{frame}
%%%%%%%%%%%%%%%%%%%%%%%%%%%%%%%%%%%%%%%%%%%%%%%%%%%%%%%%%%%%%%%%%%%%%%%%%%%%%%%%%%%%%%%%%%%%%%%%%%%%%%%%%%%%%%%%
\subsection{Operatori}
\begin{frame}[shrink=10]{Operatori}
\Minipage{.5}{	\begin{longtable}{r|l}
\text{$\setminus$*}& $\*$ \\[0.1cm]
\text{$\setminus$CF\{A\}\{B\}}& $\CF{A}{B}$ \\[0.1cm]
\text{$\setminus$codim\{A\}\{B\}}& $\codim{A}{B}$ \\[0.1cm]
\text{$\setminus$Dif\{A\}\{B\}}& $\Dif{A}{B}$ \\[0.1cm]
\text{$\setminus$Diff\{A\}\{B\}}& $\Diff{A}{B}$ \\[0.1cm]
\text{$\setminus$Floor\{A\}}& $\Floor{A}$ \\[0.1cm]
\text{$\setminus$Imm\{A\}}& $\Imm{A}$ \\[0.1cm]
\text{$\setminus$Leg\{A\}\{B\}}& $\Leg{A}{B}$ \\[0.1cm]
\text{$\setminus$nobarfrac\{A\}\{B\}}& $\nobarfrac{A}{B}$
\end{longtable}
}{.5}{\begin{longtable}{r|l}
\text{$\setminus$Norma\{A\}}& $\Norma{A}$ \\[0.1cm]
\text{$\setminus$Lie\{A\}\{B\}}& $\Lie{A}{B}$ \\[0.1cm]
\text{$\setminus$PS\{A\}\{B\}}& $\PS{A}{B}$ \\[0.1cm]
\text{$\setminus$Rad\{A\}}& $\Rad{A}$ \\[0.1cm]
\text{$\setminus$Re\{A\}}& $\Re{A}$ \\[0.1cm]
\text{$\setminus$Res\{A\}\{B\}}& $\Res{A}{B}$ \\[0.1cm]
\text{$\setminus$Stirling\{A\}\{B\}}& $\Stirling{A}{B}$ \\[0.1cm]
\text{$\setminus$Val\{A\}\{B\}}& $\Val{A}{B}$ \\[0.1cm]
\end{longtable}}
\end{frame}
%%%%%%%%%%%%%%%%%%%%%%%%%%%%%%%%%%%%%%%%%%%%%%%%%%%%%%%%%%%%%%%%%%%%%%%%%%%%%%%%%%%%%%%%%%%%%%%%%%%%%%%%%%%%%%%%
\subsection{Insiemi e gruppi}
\begin{frame}[shrink=10]{Insiemi e gruppi}
\begin{minipage}[l]{.5\textwidth}
\begin{longtable}{r|l}
\text{$\setminus$Ann\{A\}}& $\Ann{A}$ \\[0.1cm]
\text{$\setminus$AutD\{A\}}& $\AutD{A}$ \\[0.1cm]
\text{$\setminus$Aut\{A\}}& $\Aut{A}$ \\[0.1cm]
\text{$\setminus$Fix\{A\}}&$\Fix{A}$ \\[0.1cm]
\text{$\setminus$FixAut\{A\}\{B\}}& $\FixAut{A}{B}$ \\[0.1cm]
\text{$\setminus$FixAutD\{A\}\{B\}}& $\FixAutD{A}{B}$ \\[0.1cm]
\text{$\setminus$FixD\{A\}}& $\FixD{A}$ \\[0.1cm]
\text{$\setminus$FixGAut\{A\}\{B\}}& $\FixGAut{A}{B}$ \\[0.1cm]
\text{$\setminus$FixGAutD\{A\}\{B\}}& $\FixGAutD{A}{B}$ \\[0.1cm]
\text{$\setminus$GAut\{A\}}& $\GAut{A}$ \\[0.1cm]
\text{$\setminus$GAutD\{A\}}& $\GAutD{A}$ \\[0.1cm]
\text{$\setminus$GL\{A\}\{B\}}& $\GL{A}{B}$ \\[0.1cm]
\end{longtable}	
\end{minipage}\hfill
\begin{minipage}[c]{.5\textwidth}
\begin{longtable}{r|l}
\text{$\setminus$gl\{A\}}& $\gl{A}$ \\[0.1cm]
\text{$\setminus$GLE\{A\}}& $\GLE{A}$ \\[0.1cm]
\text{$\setminus$glm\{A\}\{B\}}& $\glm{A}{B}$ \\[0.1cm]
\text{$\setminus$Hom\{A\}\{B\}\{C\}}& $\Hom{A}{B}{C}$ \\[0.1cm]
\text{$\setminus$Im\{A\}}& $\Im{A}$ \\[0.1cm]
\text{$\setminus$Ker\{A\}}& $\Ker{A}$ \\[0.1cm]
\text{$\setminus$LLie\{A\}}& $\LLie{A}$ \\[0.1cm]
\text{$\setminus$OM\{A\}\{B\}}& $\OM{A}{B}$ \\[0.1cm]
\text{$\setminus$Seq\{A\}\{B\}\{C\}}& $\Seq{A}{B}{C}$ \\[0.1cm]
\text{$\setminus$Set\{A\}}& $\Set{A}$ \\[0.1cm]
\text{$\setminus$sl\{A\}}& $\sl{A}$ \\[0.1cm]
\text{$\setminus$SLM\{A\}\{B\}}& $\SLM{A}{B}$ \\[0.1cm]
\end{longtable}	
\end{minipage}
\end{frame}
%%%%%%%%%%%%%%%%%%%%%%%%%%%%%%%%%%%%%%%%%%%%%%%%%%%%%%%%%%%%%%%%%%%%%%%%%%%%%%%%%%%%%%%%%%%%%%%%%%%%%%%%%%%%%%%%
\begin{frame}{Insiemi e gruppi}
\begin{longtable}{r|l}
\text{$\setminus$slm\{A\}\{B\}}& $\slm{A}{B}$ \\[0.1cm]
\text{$\setminus$SO\{A\}\{B\}}& $\SO{A}{B}$ \\[0.1cm]
\text{$\setminus$Span\{A\}}& $\Span{A}$ \\[0.1cm]
\text{$\setminus$SU\{A\}\{B\}}& $\SU{A}{B}$ \\[0.1cm]
\text{$\setminus$Stab\{A\}\{B\}}& $\Stab{A}{B}$ \\[0.1cm]
\text{$\setminus$Sym\{A\}\{B\}}& $\Sym{A}{B}$ \\[0.1cm]
\text{$\setminus$tm\{A\}\{B\}}& $\tm{A}{B}$ \\[0.1cm]
\text{$\setminus$TM\{A\}\{B\}}& $\TM{A}{B}$ \\[0.1cm]
\text{$\setminus$Tr\{A\}}& $\Tr{A}$ \\[0.1cm]
\text{$\setminus$UM\{A\}\{B\}}& $\UM{A}{B}$ \\[0.1cm]
\end{longtable}	
\end{frame}
%%%%%%%%%%%%%%%%%%%%%%%%%%%%%%%%%%%%%%%%%%%%%%%%%%%%%%%%%%%%%%%%%%%%%%%%%%%%%%%%%%%%%%%%%%%%%%%%%%%%%%%%%%%%%%%%
\subsection{Equazioni, funzioni}
\begin{frame}[fragile, shrink=30, squeeze]{Equazioni, funzioni}
\begin{verbatim}\EqL{A}{B}\end{verbatim}
\EqL{Equazione\ con\ label}{B}
\begin{verbatim}\Eq{A & B\\ C & D}\end{verbatim}
\Eq{A & B\\ C & D}
\begin{verbatim}\Map{A}{B}{C}{D}{E}\end{verbatim}
$$\Map{A}{B}{C}{D}{E}$$
\begin{verbatim}\Map{A}{B}{C}{D}\end{verbatim}
$$\NMap{A}{B}{C}{D}$$
\end{frame}
%%%%%%%%%%%%%%%%%%%%%%%%%%%%%%%%%%%%%%%%%%%%%%%%%%%%%%%%%%%%%%%%%%%%%%%%%%%%%%%%%%%%%%%%%%%%%%%%%%%%%%%%%%%%%%%%
\subsection{Box, minipage, quotes}
\begin{frame}[fragile, shrink=10]{Box, minipage, quotes}
\begin{verbatim}
\begin{TitoloIntro}[colbacktitle=red, width=\textwidth]{Titolo}
{Riquadro}\end{TitoloIntro}	
\end{verbatim}
\begin{TitoloIntro}[colbacktitle=red, width=\textwidth]{Titolo}{Riquadro}\end{TitoloIntro}
\begin{verbatim}\Minipage{0.45}{Pagina a sinistra}{0.45}{Pagina a destra}\end{verbatim}
\Minipage{0.45}{Pagina a sinistra\newline\newline\newline\newline {\color{white}.} }{0.45}{Pagina a destra\newline\newline \Quote{A}{B}{.5}}	
\vspace{-2cm}
\begin{verbatim}\Quote{A}{B}{.5}\end{verbatim}
\end{frame}
%%%%%%%%%%%%%%%%%%%%%%%%%%%%%%%%%%%%%%%%%%%%%%%%%%%%%%%%%%%%%%%%%%%%%%%%%%%%%%%%%%%%%%%%%%%%%%%%%%%%%%%%%%%%%%%%
\subsection{Altre funzionalità}
\begin{frame}[fragile]{Altre funzionalità}
Nel file custom sono implementati altri pacchetti che supportano diversi comandi e funzionalità utili:
\begin{verbatim}polynomial
accents
tikz
pgf\end{verbatim}
Chi vuole approfondire può guardare ad esempio ai template book, article al seguente link:\newline
\url{https://poisson.phc.dm.unipi.it/~puddu/Templates/Templates.html}
\end{frame}
%%%%%%%%%%%%%%%%%%%%%%%%%%%%%%%%%%%%%%%%%%%%%%%%%%%%%%%%%%%%%%%%%%%%%%%%%%%%%%%%%%%%%%%%%%%%%%%%%%%%%%%%%%%%%%%%
\subsubsection{Esempio di utilizzo tikz3d-plot}
\begin{frame}{Esempio di utilizzo tikz3d-plot}
	\tdplotsetmaincoords{60}{110}
	\begin{center}
    \begin{tikzpicture}[scale=4.5, tdplot_main_coords]
    \pgfmathsetmacro{\rvec}{.8}
    \pgfmathsetmacro{\thetavec}{30}
    \pgfmathsetmacro{\phivec}{60}
    \coordinate (O) at (0,0,0);
    \draw[thick,->] (0,0,0) -- (1,0,0) node[anchor=north east]{$x$};
    \draw[thick,->] (0,0,0) -- (0,1,0) node[anchor=north west]{$y$};
    \draw[thick,->] (0,0,0) -- (0,0,1) node[anchor=south]{$z$};
    \tdplotsetcoord{P}{\rvec}{\thetavec}{\phivec}
    \draw[-stealth,color=red] (O) -- (P);
    \draw[dashed, color=red] (O) -- (Pxy);
    \draw[dashed, color=red] (P) -- (Pxy);
    \tdplotdrawarc{(O)}{0.2}{0}{\phivec}{anchor=north}{$\phi$}
    \tdplotsetthetaplanecoords{\phivec}
    \tdplotdrawarc[tdplot_rotated_coords]{(0,0,0)}{0.5}{0}%
        {\thetavec}{anchor=south west}{$\theta$}
    \draw[dashed,tdplot_rotated_coords] (\rvec,0,0) arc (0:90:\rvec);
    \draw[dashed] (\rvec,0,0) arc (0:90:\rvec);
    \tdplotsetrotatedcoords{\phivec}{\thetavec}{0}
    \tdplotsetrotatedcoordsorigin{(P)}
    \draw[thick,tdplot_rotated_coords,->] (0,0,0)
        -- (.5,0,0) node[anchor=north west]{$x’$};
    \draw[thick,tdplot_rotated_coords,->] (0,0,0)
        -- (0,.5,0) node[anchor=west]{$y’$};
    \draw[thick,tdplot_rotated_coords,->] (0,0,0)
        -- (0,0,.5) node[anchor=south]{$z’$};
    \draw[-stealth,color=blue,tdplot_rotated_coords] (0,0,0) -- (.2,.2,.2);
    \draw[dashed,color=blue,tdplot_rotated_coords] (0,0,0) -- (.2,.2,0);
    \draw[dashed,color=blue,tdplot_rotated_coords] (.2,.2,0) -- (.2,.2,.2);
    \tdplotdrawarc[tdplot_rotated_coords,color=blue]{(0,0,0)}{0.2}{0}%
        {45}{anchor=north west,color=black}{$\phi’$}
    \tdplotsetrotatedthetaplanecoords{45}
    \tdplotdrawarc[tdplot_rotated_coords,color=blue]{(0,0,0)}{0.2}{0}%
        {55}{anchor=south west,color=black}{$\theta’$}
\end{tikzpicture}
	\end{center}
\end{frame}
%%%%%%%%%%%%%%%%%%%%%%%%%%%%%%%%%%%%%%%%%%%%%%%%%%%%%%%%%%%%%%%%%%%%%%%%%%%%%%%%%%%%%%%%%%%%%%%%%%%%%%%%%%%%%%%%
\subsubsection{Esempio di utilizzo pgfplots}
\begin{frame}{Esempio di utilizzo pgfplots}
\begin{center}
\begin{tikzpicture}
\begin{axis}[
    colormap/PuBu,
    samples=12, domain=-5:5, y domain=-4:4,  xmin=-5, xmax=5, ymin=-4, ymax=4]
    \addplot3 [surf,     samples=24,    samples y=40] {x^2+y^2-1};
\end{axis}
\end{tikzpicture}	
\end{center}
\end{frame}
%%%%%%%%%%%%%%%%%%%%%%%%%%%%%%%%%%%%%%%%%%%%%%%%%%%%%%%%%%%%%%%%%%%%%%%%%%%%%%%%%%%%%%%%%%%%%%%%%%%%%%%%%%%%%%%%
%%%%%%%%%%%%%%%%%%%%%%%%%%%%%%%%%%%%%%%%%%%%%%%%%%%%%%%%%%%%%%%%%%%%%%%%%%%%%%%%%%%%%%%%%%%%%%%%%%%%%%%%%%%%%%%%
\part{Bibliografia}
\begin{frame}{Bibliografia}{}
\begin{thebibliography}{}
\setbeamertemplate{bibliography item}[online]
%Opzioni disponibili: default|article|book|online|triangle|text

        \bibitem[Puddu, 2021]{Puddu}
        M.~Puddu.
        \newblock {\em Sito Web}.
        \newblock {\em \url{poisson.phc.dm.unipi.it/~puddu}}

    \setbeamertemplate{bibliography item}[text]
        \bibitem[Autore, 2021]{AutoreLabel}
        N.~Cognome.
        \newblock {\em Titolo}.
        \newblock Publisher, Anno.

\end{thebibliography}
\end{frame}
%%%%%%%%%%%%%%%%%%%%%%%%%%%%%%%%%%%%%%%%%%%%%%%%%%%%%%%%%%%%%%%%%%%%%%%%%%%%%%%%%%%%%%%%%%%%%%%%%%%%%%%%%%%%%%%%
{\setbeamercolor{background canvas}{bg=MyBlue!80}
\begin{frame}[plain]
    \begin{center}\bf\Huge\color{white}
    Grazie per l'attenzione!
    \end{center}
\end{frame}}
%%%%%%%%%%%%%%%%%%%%%%%%%%%%%%%%%%%%%%%%%%%%%%%%%%%%%%%%%%%%%%%%%%%%%%%%%%%%%%%%%%%%%%%%%%%%%%%%%%%%%%%%%%%%%%%%
\end{document}

\end{document}